% \documentclass[conference]{IEEEtran}
\documentclass[sustainability,article,submit,moreauthors,pdftex]{Definitions/mdpi}

\firstpage{1}
\makeatletter
\setcounter{page}{\@firstpage}
\makeatother
\pubvolume{1}
\issuenum{1}
\articlenumber{0}
\pubyear{2026}
\copyrightyear{2026}
\externaleditor{~} % More than 1 editor, please add `` and '' before the last editor name
\datereceived{  }
\daterevised{ } % Comment out if no revised date
\dateaccepted{ }
\datepublished{ }

% \usepackage{amsmath,amssymb,amsfonts}
% \usepackage{graphicx}
% \usepackage{booktabs}
% \usepackage{multirow}
% \usepackage{array}
% \usepackage{float}
% \usepackage{subcaption}
% \usepackage{tikz}
% \usepackage{pgfplots}
% \usepackage{algorithm}
% \usepackage{algpseudocode}
% \usepackage[hidelinks]{hyperref}
% \usepackage{caption}

\usepackage{adjustbox}
\usepackage{booktabs}
\usepackage{makecell}
\usepackage{multirow}
\usepackage{rotating}
\usepackage{array}
\usepackage{xcolor}
\usepackage{colortbl}
\usepackage{pifont}
\usepackage{tikz}
\usetikzlibrary{shapes.geometric, arrows.meta, positioning}
\captionsetup[table]{justification=centering}
% \pgfplotsset{compat=1.18}

% Float placement: allow more/larger floats per page so figures stay near
% their text instead of being deferred to the end of the document.
\usepackage{placeins}
\renewcommand{\topfraction}{0.9}
\renewcommand{\bottomfraction}{0.9}
\renewcommand{\textfraction}{0.08}
\renewcommand{\floatpagefraction}{0.7}
\setcounter{topnumber}{3}
\setcounter{bottomnumber}{2}
\setcounter{totalnumber}{5}

\newcommand{\cmark}{\ding{51}}   % checkmark
\newcommand{\pmark}{$\circ$}     % partial
\newcommand{\xmark}{---}         % not included

% Highlight color for "This Paper" column
\definecolor{thispaper}{RGB}{232,244,253}

\Title{A Locational Price and Carbon Intensity-Aware \\Charger Allocation and Charge--Discharge Scheduling Strategy for Electric Vehicle Fleets}

% \TitleCitation{Title}
% \author{
% Jiwan Jung
% jung232@wisc.edu
% Univeristy of Wisconsin-Madison
% }

\newcommand{\orcidauthorA}{0000-0003-3662-335X}

\Author{Jiwan Jung $^{1}$, Shitikantha Dash $^{2}$\orcidA{} and Dipti Srinivasan $^{2,}$*}

\AuthorNames{Jiwan Jung, Shitikantha Dash, Dipti Srinivasan}

\address{%
$^{1}$ \quad Department of Electrical and Computer Engineering, University of Wisconsin-Madison, \linebreak Madison, USA; {jung232@wisc.edu}\\
$^{2}$ \quad Department of Electrical and Computer Engineering, National University of Singapore, \linebreak Singapore 117538, Singapore; {sdash@nus.edu.sg} }

\corres{Correspondence: {dipti@nus.edu.sg} }

% \maketitle
% =========================================================

\abstract{%
The rapid growth of electric vehicle adoption introduces new challenges for grid operators, as uncoordinated charging concentrates demand during stressed periods and ignores both economic and environmental signals. This paper proposes a multi-objective EV charging and discharging scheduling framework that combines locational marginal prices and nodal carbon intensity into a single integrated charging signal. The framework uses a two-stage pipeline: a capacity-aware charging allocation procedure that pre-allocates EVs to feasible charging nodes, followed by a linear program that optimizes charge and discharge schedules for each EV. Vehicle-to-grid discharging is included with a linear battery degradation cost, and a soft carbon cap is imposed relative to an uncoordinated baseline. Three scheduling schemes are evaluated, and a sweep over the cost--emissions weighting reveals the full trade-off frontier. Cost is reported as total operating cost, comprising energy cost and battery degradation, which is the quantity the scheduler minimizes. The balanced scheme reduces total operating cost by 28.4\% and carbon emissions by 21.6\% relative to the uncoordinated baseline, showing that both objectives can be improved substantially and simultaneously.}

% =========================================================

\keyword{Electric vehicles, vehicle-to-grid, optimization, carbon-aware, cost-aware, locational marginal pricing, linear programming, smart grid}

% =========================================================

\begin{document}

\section*{\textcolor{red}{Nomenclature}}

\noindent\textit{A. Sets and Indices}\vspace{6pt}\\

\begin{tabular}{@{} p{2.2cm} p{10.3cm} }
$\mathcal{E}$ & Set of all EVs \\
$\mathcal{E}_n$ & Set of EVs assigned to node $n$ \\
$\mathcal{E}_{n,t}$ & Set of EVs assigned to node $n$ and active in period $t$ \\
$\mathcal{N}$ & Set of charging nodes \\
$\mathcal{T}$ & Set of discrete time periods \\
$\mathcal{T}_i$ & Charging period of EV $i$ \\
$\mathcal{T}^{\mathrm{base}}_i$ & Periods in which EV $i$ charges under the uncoordinated baseline \\
\end{tabular}\vspace{10pt}

\noindent\textit{B. Decision Variables}\vspace{6pt}\\

\begin{tabular}{@{} p{2.2cm} p{10.3cm} }
$P^{\mathrm{c}}_{i,t}$ & Charging power of EV $i$ at time $t$ (kW) \\
$P^{\mathrm{d}}_{i,t}$ & Discharging power of EV $i$ at time $t$ (kW) \\
$E_{i,t}$ & Battery energy state of EV $i$ at time $t$ (kWh) \\
$s^{\mathrm{carbon}}$ & Soft slack for carbon cap constraint (g CO$_2$) \\
\end{tabular}\vspace{10pt}

\noindent\textit{B1. Stage-1 Output (fixed before the LP is solved)}\vspace{4pt}\\

\begin{tabular}{@{} p{2.2cm} p{10.3cm} }
$n_i$ & Node assigned to EV $i$ \\
\end{tabular}\vspace{10pt}

\noindent\textit{C. Parameters}\vspace{6pt}\\

\noindent\textit{C1. Charging Allocation}\vspace{4pt}\\

\begin{tabular}{@{} p{2.0cm} p{9cm} p{1.5cm} }
$t_i^{\mathrm{arr}}$ & Arrival time of EV $i$ & h \\
$t_i^{\mathrm{dep}}$ & Departure time of EV $i$ & h \\
$E_i^{\mathrm{arr}}$ & Arrival energy of EV $i$ & kWh \\
$E_i^{\mathrm{dep}}$ & Required departure energy of EV $i$ & kWh \\
$E_i^{\mathrm{cap}}$ & Battery capacity of EV $i$ & kWh \\
$P_i^{\mathrm{c,max}}$ & Maximum charging power of EV $i$ & kW \\
$P_i^{\mathrm{d,max}}$ & Maximum discharging power of EV $i$ & kW \\
$\eta_c$ & Charging efficiency & \\
$\eta_d$ & Discharging efficiency & \\
$\delta$ & Battery degradation cost coefficient & \$/kWh \\
$P_n^{\max}$ & Maximum power capacity of node $n$ & kW \\
\end{tabular}\vspace{10pt}

\noindent\textit{C2. Nodal Signals \& State}\vspace{4pt}\\

\begin{tabular}{@{} p{2.0cm} p{9cm} p{1.5cm} }
$p_{n,t}$ & Locational marginal price at node $n$, period $t$ & \$/kWh \\
$\hat{p}_{n,t}$ & Normalized locational marginal price & \\
$p_{\max}$ & Maximum LMP over all nodes and periods & \$/kWh \\
$p_{\min}$ & Minimum LMP over all nodes and periods & \$/kWh \\
$c_{n,t}$ & Nodal carbon intensity at node $n$, period $t$ & g CO$_2$/kWh \\
$\hat{c}_{n,t}$ & Normalized carbon intensity & \\
$c_{\max}$ & Maximum carbon intensity over all nodes and periods & g CO$_2$/kWh \\
$S_{n,t}$ & Integrated charging signal at node $n$, period $t$ & \\
\end{tabular}\vspace{10pt}

\noindent\textit{C3. Assignment Heuristic}\vspace{4pt}\\

\begin{tabular}{@{} p{2.2cm} p{10.3cm} }
$Score_{i,n}$ & Score given to node $n$ for EV $i$ \\
$d_{i,n}$ & Euclidean distance between EV $i$ and node $n$ (grid units) \\
$d_{\max}$ & Maximum allowable assignment distance (grid units) \\
\end{tabular}\vspace{10pt}

\noindent\textit{C4. Weighting Coefficients}\vspace{4pt}\\

\begin{tabular}{@{} p{2.2cm} p{10.3cm} }
$\alpha$ & Cost weighting coefficient \\
$\beta$ & Carbon weighting coefficient \\
$\phi(\beta)$ & Effective carbon cap fraction at weight $\beta$ \\
$\phi_0$ & Tightest allowable carbon cap fraction \\
$M$ & Penalty coefficient for soft carbon constraint \\
\end{tabular}\newpage

\section{Introduction}

\subsection{Background and Motivation}

The global adoption of electric vehicles (EVs) continues to grow as governments promote them as cleaner alternatives to internal combustion engines \cite{EV_adoption_growth}. The transition also creates operational problems: a large fleet charging on arrival concentrates demand into a few hours, raising peak loads and stressing distribution feeders and voltage stability \cite{uncoordinated_charging_grid_impacts}.

Timing carries economic weight as well. Locational marginal prices vary both over time and across nodes, so uncoordinated charging exposes owners to avoidable cost.

The environmental dimension is equally dynamic. Grid carbon intensity tracks the generation mix and varies by location as well as by hour, so the same energy carries a different carbon footprint depending on when and where it is drawn. Ignoring both signals forgoes the chance to align charging with supply that is simultaneously cheap and clean.

Bidirectional charging extends the opportunity further. Vehicle-to-grid (V2G) operation lets EVs discharge into the grid when price or carbon intensity is high, acting as distributed storage that earns revenue and supports grid operation.

\subsection{Previous Works and Gaps}

Table~\ref{tab:literature_review} compares this work against representative
studies on EV charging and discharging scheduling. Three gaps emerge.

\textit{Single-signal scheduling.} Most schedulers optimize one signal,
either cost~\cite{sortomme2011, sundstrom2012, gan2013, li2014dlmp,
huang2015dlmp, sarker2016} or emissions~\cite{hoehne2016, xu2020ghg,
cheng2022carbon}. The multi-objective formulations that treat
both~\cite{das2020, aurangzeb2025, alikhani2025, ambrosino2025} use
time-of-use or zonal prices and system-wide emission factors, so neither
signal is resolved by location and the inter-nodal differences this work
exploits remain invisible to them.

\textit{Decoupled node assignment.} Where EVs charge is usually decided
separately from when they charge, if at all. Only~\cite{moradipari2020,
arias2023lmp} assign EVs to charging locations, and they do so through
pricing or queueing rather than by checking feasibility. No prior work
verifies, before scheduling begins, that an assignment jointly satisfies
distance, nodal power capacity, and energy-throughput limits, so the
resulting schedule may be physically unrealizable.

\textit{Unconstrained emissions with V2G.} Bidirectional
schedulers~\cite{sarker2016, hoehne2016, das2020, aurangzeb2025,
alikhani2025} rarely bound fleet emissions, and~\cite{hoehne2016} shows
that mistimed V2G can raise them outright. None prices battery
degradation while enforcing a carbon cap.

No prior work combines nodal LMP and nodal carbon intensity in one
scheduling signal, embeds it in a feasibility-guaranteed allocation
stage, and optimizes V2G under degradation costs and a carbon cap within
a single linear program.

\begin{table*}[t]
\renewcommand{\arraystretch}{1.15}
\caption{Comparison of Related Work on EV Charging and Discharging Scheduling}\label{tab:literature_review}
\begin{adjustwidth}{-\extralength}{0cm}\setlength{\tabcolsep}{2mm}
\centering
\resizebox{\linewidth}{!}{%
\scriptsize
\setlength{\tabcolsep}{3.5pt}
\begin{tabular}{
  l
  >{\centering\arraybackslash}p{0.95cm}
  >{\centering\arraybackslash}p{0.95cm}
  >{\centering\arraybackslash}p{0.95cm}
  >{\centering\arraybackslash}p{0.95cm}
  >{\centering\arraybackslash}p{0.95cm}
  >{\centering\arraybackslash}p{0.95cm}
  >{\centering\arraybackslash}p{0.95cm}
  >{\centering\arraybackslash}p{0.95cm}
  >{\centering\arraybackslash}p{0.95cm}
  >{\centering\arraybackslash}p{0.95cm}
  >{\centering\arraybackslash}p{0.95cm}
  >{\centering\arraybackslash}p{0.95cm}
  >{\columncolor{thispaper}\centering\arraybackslash}p{1.05cm}
}
\toprule
\multirow{2}{*}{\textbf{Feature}} &
  \rotatebox{90}{\parbox{2.6cm}{Sortomme\\et al.~\cite{sortomme2011} (2011)}} &
  \rotatebox{90}{\parbox{2.6cm}{Sundstr\"om \&\\Binding~\cite{sundstrom2012} (2012)}} &
  \rotatebox{90}{\parbox{2.6cm}{Gan\\et al.~\cite{gan2013} (2013)}} &
  \rotatebox{90}{\parbox{2.6cm}{Li\\et al.~\cite{li2014dlmp} (2014)}} &
  \rotatebox{90}{\parbox{2.6cm}{Sarker\\et al.~\cite{sarker2016} (2016)}} &
  \rotatebox{90}{\parbox{2.6cm}{Hoehne \&\\Chester~\cite{hoehne2016} (2016)}} &
  \rotatebox{90}{\parbox{2.6cm}{Das\\et al.~\cite{das2020} (2020)}} &
  \rotatebox{90}{\parbox{2.6cm}{Moradipari \&\\Alizadeh~\cite{moradipari2020} (2020)}} &
  \rotatebox{90}{\parbox{2.6cm}{Cheng\\et al.~\cite{cheng2022carbon} (2022)}} &
  \rotatebox{90}{\parbox{2.6cm}{Arias\\et al.~\cite{arias2023lmp} (2023)}} &
  \rotatebox{90}{\parbox{2.6cm}{Alikhani\\et al.~\cite{alikhani2025} (2025)}} &
  \rotatebox{90}{\parbox{2.6cm}{Ambrosino\\et al.~\cite{ambrosino2025} (2025)}} &
  \rotatebox{90}{\parbox{2.6cm}{\textbf{This Paper}\\(2026)}} \\[0.2cm]
\midrule
\multicolumn{14}{l}{\textit{Economic Signal}} \\
\quad Wholesale market price signal
  & \xmark & \cmark & \xmark & \cmark & \cmark & \xmark & \xmark & \xmark & \xmark & \cmark & \cmark & \cmark & \cellcolor{thispaper}\cmark \\
\quad Locational (nodal) prices (LMP/DLMP)
  & \xmark & \xmark & \xmark & \cmark & \xmark & \xmark & \xmark & \xmark & \xmark & \cmark & \xmark & \xmark & \cellcolor{thispaper}\cmark \\
\midrule
\multicolumn{14}{l}{\textit{Environmental Signal}} \\
\quad Carbon-aware objective
  & \xmark & \xmark & \xmark & \xmark & \xmark & \cmark & \cmark & \xmark & \cmark & \xmark & \cmark & \cmark & \cellcolor{thispaper}\cmark \\
\quad Time-varying carbon intensity
  & \xmark & \xmark & \xmark & \xmark & \xmark & \cmark & \cmark & \xmark & \cmark & \xmark & \cmark & \cmark & \cellcolor{thispaper}\cmark \\
\quad Nodal carbon intensity
  & \xmark & \xmark & \xmark & \xmark & \xmark & \pmark & \xmark & \xmark & \xmark & \xmark & \pmark & \xmark & \cellcolor{thispaper}\cmark \\
\midrule
\multicolumn{14}{l}{\textit{Multi-Objective Formulation}} \\
\quad Cost \& carbon jointly optimized
  & \xmark & \xmark & \xmark & \xmark & \xmark & \xmark & \cmark & \xmark & \xmark & \xmark & \cmark & \cmark & \cellcolor{thispaper}\cmark \\
\quad Cost--carbon trade-off (Pareto) analysis
  & \xmark & \xmark & \xmark & \xmark & \xmark & \xmark & \pmark & \xmark & \xmark & \xmark & \cmark & \pmark & \cellcolor{thispaper}\cmark \\
\quad Carbon-cap constraint
  & \xmark & \xmark & \xmark & \xmark & \xmark & \xmark & \xmark & \xmark & \xmark & \xmark & \xmark & \xmark & \cellcolor{thispaper}\pmark \\
\midrule
\multicolumn{14}{l}{\textit{V2G / Discharging}} \\
\quad Vehicle-to-grid (V2G) discharging
  & \xmark & \xmark & \xmark & \xmark & \cmark & \cmark & \cmark & \xmark & \xmark & \xmark & \cmark & \xmark & \cellcolor{thispaper}\cmark \\
\quad Battery degradation cost modeled
  & \xmark & \xmark & \xmark & \xmark & \cmark & \xmark & \cmark & \xmark & \xmark & \xmark & \xmark & \xmark & \cellcolor{thispaper}\cmark \\
\midrule
\multicolumn{14}{l}{\textit{Node Assignment}} \\
\quad EV-to-node assignment modeled
  & \xmark & \xmark & \xmark & \xmark & \xmark & \xmark & \xmark & \cmark & \xmark & \cmark & \xmark & \xmark & \cellcolor{thispaper}\cmark \\
\quad Nodal power capacity constraints
  & \xmark & \cmark & \xmark & \cmark & \xmark & \xmark & \xmark & \pmark & \cmark & \cmark & \xmark & \pmark & \cellcolor{thispaper}\cmark \\
\quad Distance-feasibility constraint
  & \xmark & \xmark & \xmark & \xmark & \xmark & \xmark & \xmark & \pmark & \xmark & \xmark & \xmark & \xmark & \cellcolor{thispaper}\cmark \\
\midrule
\multicolumn{14}{l}{\textit{Scheduling Framework}} \\
\quad Integrated price+carbon signal (ICS)
  & \xmark & \xmark & \xmark & \xmark & \xmark & \xmark & \xmark & \xmark & \xmark & \xmark & \xmark & \xmark & \cellcolor{thispaper}\cmark \\
\quad Solution method
  & QP & LP & Dec.\ conv. & DLMP & MILP & Scen. & MOO & Q+O & LP & MISOCP & MOO & LP & \cellcolor{thispaper}LP \\
\bottomrule
\end{tabular}}
\end{adjustwidth}

\vspace{2pt}

{\footnotesize
\cmark~= included;\quad
\pmark~= partial (feature present but limited in scope);\quad
\xmark~= not considered.\\[2pt]
\textit{Partial} examples: Hoehne \& Chester use regional rather than nodal
emission profiles; Alikhani et al.~use zonal marginal emission factors;
Das et al.~compute Pareto solutions over four objectives rather than a
dedicated cost--carbon frontier; Ambrosino et al.~sweep three discrete
cost--carbon weights and assume a non-binding station capacity;
Moradipari \& Alizadeh capture station congestion and en-route station
choice rather than explicit nodal power caps and distance limits. The
carbon cap in this work is marked partial because it is formulated and
enforced but, under the $\beta$-coupled cap schedule adopted here, does
not bind at any reported operating point (Section~\ref{sec:overall}).\\[2pt]
Methods: QP = quadratic programming; LP = linear programming;
Dec.\ conv.\ = decentralized convex optimization; DLMP = DLMP-based
congestion pricing; MILP = mixed-integer LP; Scen.\ = scenario-based
emissions scheduling; MOO = multi-objective (evolutionary/metaheuristic)
optimization; Q+O = queueing model with optimization; MISOCP =
mixed-integer second-order cone programming.
}
\end{table*}

\subsection{Contributions}

Based on the discussion above, the key contributions of this paper can be listed as follows:

\begin{enumerate}

\item \textit{Integrated Charging Signal (ICS)} that combines locational marginal pricing (LMP) and nodal carbon intensity into a unified scheduling signal for simultaneous consideration of economic and environmental objectives.

\item \textit{Capacity-aware Charging Allocation Framework} that pre-assigns EVs to charging nodes while accounting for distance constraints, nodal power limits, and energy-throughput feasibility.

\item \textit{Linear programming based charging and discharging scheduler} that optimizes charging and discharging decisions using the proposed ICS while incorporating battery degradation costs and carbon-cap constraints.

\end{enumerate}

% =========================================================

\section{Proposed Multi-Objective EV Charging Framework}

\subsection{Framework Overview}

\begin{figure}[tbp]
\centering
\begin{tikzpicture}[
  font=\footnotesize,
  node distance=3mm,
  arrow/.style={-{Latex[length=2mm]}, thick},
  term/.style={draw, ellipse, minimum width=2.2cm, minimum height=0.7cm,
               align=center},
  io/.style={draw, trapezium, trapezium left angle=70, trapezium right angle=110,
             trapezium stretches body, text width=7cm, align=left, inner sep=4pt},
  proc/.style={draw, rectangle, text width=7.4cm, align=left, inner sep=5pt}
]
\node[term] (start) {\textbf{Start}};
\node[io, below=of start] (input) {%
  {\centering\textbf{Inputs}\par}\vspace{0.3mm}
  1) EV data: arrival/departure, location, initial SOC\\
  2) Grid data: nodal LMP $p_{n,t}$, carbon intensity $c_{n,t}$};
\node[proc, below=of input] (stage1) {%
  {\centering\textbf{Stage 1: Charging Allocation}\par}\vspace{0.3mm}
  ICS \eqref{eq:composite_signal}, node score \eqref{eq:assignment_score}\\
  s.t.\ (1) distance feasibility \eqref{eq:distance_constraint}\\
  \phantom{s.t.\ }(2) peak-power feasibility \eqref{eq:peak_power}\\
  \phantom{s.t.\ }(3) energy-throughput feasibility \eqref{eq:energy_throughput}\\[0.3mm]
  {\centering\textit{greedy assignment procedure}\par}};
\node[proc, below=of stage1] (stage2) {%
  {\centering\textbf{Stage 2: LP Scheduling}\par}\vspace{0.3mm}
  $\min$ \eqref{eq:objective}\\
  s.t.\ (1) battery dynamics \eqref{eq:battery_dynamics}, SOC limits \eqref{eq:init_soc}--\eqref{eq:battery_limit}\\
  \phantom{s.t.\ }(2) charge/discharge limits \eqref{eq:charge_limit}, \eqref{eq:discharge_limit}\\
  \phantom{s.t.\ }(3) node capacity limits \eqref{eq:node_limit}\\
  \phantom{s.t.\ }(4) carbon cap \eqref{eq:carbon_cap}};
\node[io, below=of stage2] (output) {%
  {\centering\textbf{Output}\par}\vspace{0.3mm}
  Return $\{P^{\mathrm{c}}_{i,t}, P^{\mathrm{d}}_{i,t}\}$, total cost, total emissions};
\node[term, below=of output] (stop) {\textbf{Stop}};
\draw[arrow] (start)  -- (input);
\draw[arrow] (input)  -- (stage1);
\draw[arrow] (stage1) -- (stage2);
\draw[arrow] (stage2) -- (output);
\draw[arrow] (output) -- (stop);
\end{tikzpicture}
\caption{Overview of the proposed two-stage EV charging framework.}
\label{fig:framework}
\end{figure}

As depicted in Figure~\ref{fig:framework}, the framework consists of the following components.

First, EV data---including arrival and departure times, location, and state of charge---is provided as input. Grid data, comprising the LMP signal and carbon intensity, is also provided to the scheduler. The framework then follows a two-stage architecture to schedule and optimize fleet charging.

\textbf{Stage 1: Charging Allocation.} Each EV is pre-assigned to a charging node using a capacity-aware greedy procedure driven by the Integrated Charging Signal (ICS). The procedure enforces distance, peak-power, and energy-throughput constraints to guarantee that every assigned EV can be fully charged within its departure window. The full formulation is given in Section~\ref{sec:formulation}.

\textbf{Stage 2: LP Scheduling.} With the node assignments produced by Stage 1 held fixed, a linear program optimizes per-EV charge and discharge profiles across the 24-hour scheduling horizon. The LP takes LMP prices, grid carbon intensity, and EV arrival and departure information as inputs and produces per-EV schedules.

Both stages are driven by the same weighting $(\alpha, \beta)$, and the assignments are fixed only within a single run. Changing the weighting re-scores the ICS in~\eqref{eq:assignment_score} and therefore changes which node each EV is assigned to, before the LP is solved. Every operating point reported in this paper is the result of running both stages at that point's weighting.

Finally, total cost and carbon emissions are computed from the resulting schedule.

\subsection{EV System Representation}

Each EV is represented by a tuple of parameters: arrival time, departure time, arrival energy $E_i^{\mathrm{arr}}$, required departure energy $E_i^{\mathrm{dep}}$, battery capacity $E_i^{\mathrm{cap}}$, maximum charge and discharge power limits, and spatial grid coordinates depicting the starting location $(x_i, y_i)$. The arrival and departure times of each EV are treated as fixed, known inputs.

\subsection{Integrated Charging Signal}

This study introduces the Integrated Charging Signal (ICS), a single scalar signal composed of the LMP price and carbon intensity, capturing the economic and environmental cost of charging at node $n$ in period $t$.

When the ICS is low, electricity is cheap and the grid is relatively clean, so the scheduler is incentivized to charge. When the ICS is high, electricity is expensive and carbon-intensive, so the scheduler will delay charging or discharge back to the grid to reduce cost.

The weighting parameters $\alpha$ and $\beta$ allow the scheduler to trade off between cost and carbon emission objectives. The formal derivation, including normalization, is given in Section~\ref{sec:formulation}.

\section{Problem Formulation}
\label{sec:formulation}

\subsection{Signal Normalization}

Electricity price and carbon intensity operate on different scales and units, requiring normalization before they can be combined into a single scheduling signal.

Normalization is performed relative to the maximum observed value:

\begin{equation}
\hat{p}_{n,t}
=
\frac{p_{n,t}}{p_{\max}}
\label{eq:price_norm}
\end{equation}

\begin{equation}
\hat{c}_{n,t}
=
\frac{c_{n,t}}{c_{\max}}
\label{eq:carbon_norm}
\end{equation}

where $p_{\max}$ and $c_{\max}$ are the maximum values of the LMP signal and nodal carbon intensity over the 24-hour period.

Scaling by the maximum is preferred here to a min--max rescaling. Min--max normalization maps the lowest observed price to exactly zero, so the scheduler would treat charging in that period as costless, and it would discard the information carried by negative LMPs, where consumption is actively rewarded. Dividing by the maximum preserves both the sign and the ratio structure of the original signals: $\hat{c}_{n,t}$ lies in $[0,1]$, and $\hat{p}_{n,t}$ lies in $[p_{\min}/p_{\max},\,1]$, which is negative whenever prices fall below zero.

The two normalized signals are then combined into a single Integrated Charging Signal, $S_{n,t}$:

\begin{equation}
    S_{n,t}
    =
    \alpha\,\hat{p}_{n,t}
    +
    \beta\,\hat{c}_{n,t}
    \quad \forall\ n \in \mathcal{N},\ t \in \mathcal{T}
    \label{eq:composite_signal}
\end{equation}

subject to:

\begin{equation}
    \alpha + \beta = 1,
    \qquad
    \alpha,\,\beta \in [0,1]
    \label{eq:weights}
\end{equation}

The weight $\alpha$ controls preference given to electricity cost minimization, while $\beta$ controls preference given to carbon emission reduction.

Because each signal is scaled by its own maximum, equal weights do not imply equal influence on the schedule. Over the study period $\hat{p}_{n,t}$ spans $[-0.314,\,1]$ while $\hat{c}_{n,t}$ spans $[0.594,\,1]$, so the price term carries roughly 3.2 times the temporal variation of the carbon term. The scheduler responds only to variation in $S_{n,t}$ over time, so the balanced scheme lies closer to the cost-only endpoint than its label suggests, which is consistent with its position on the frontier reported in Section~\ref{sec:pareto}.

\subsection{Charging Allocation Formulation}

\subsubsection{Feasible Charging Allocation}

Each EV is assigned to exactly one charging node, denoted $n_i \in \mathcal{N}$.

Each node is scored using the cumulative ICS values over EV $i$'s active window, $\mathcal{T}_i$:

\begin{equation}
Score_{i,n}
=
\sum_{t \in \mathcal{T}_i} S_{n,t}
\label{eq:assignment_score}
\end{equation}

Once the node with the minimum score is identified, distance feasibility is enforced to ensure the EV is not assigned to a node beyond $d_{\max}$:

\begin{equation}
d_{i,n} \le d_{\max}
\label{eq:distance_constraint}
\end{equation}

where $d_{i,n}$ is the Euclidean distance between EV $i$ and node $n$. If no node satisfies~\eqref{eq:distance_constraint} for a given EV, the constraint is relaxed for that vehicle and a fallback assignment is applied.

\subsubsection{Greedy Assignment}

The assignment problem is solved greedily rather than globally, in two passes. First, all feasible $(i,n)$ pairs satisfying~\eqref{eq:distance_constraint} are ranked in ascending order of $Score_{i,n}$. An assignment is accepted only if two conditions hold: \eqref{eq:peak_power} must hold at every $t \in \mathcal{T}_i$, and \eqref{eq:energy_throughput} must hold over the window as a whole.

\begin{equation}
    \sum_{j \in \mathcal{E}_{n,t}} P^{\mathrm{c,max}}_{j}
    \;+\;
    P^{\mathrm{c,max}}_{i}
    \;\le\;
    P_n^{\max}
    \label{eq:peak_power}
\end{equation}

where $\mathcal{E}_{n,t}$ is the set of EVs already committed to node $n$ that are active in period $t$. Equation~\eqref{eq:peak_power} checks that adding EV $i$ does not push the total maximum charging power at node $n$ above its capacity limit in any period.

The second condition ensures that energy-throughput is feasible for the EV:

\begin{equation}
    \sum_{t \in \mathcal{T}_i}
    \Bigl( P_n^{\max} - \!\!\sum_{j \in \mathcal{E}_{n,t}} \!\! P^{\mathrm{c,max}}_{j} \Bigr)
    \;\ge\;
    \frac{E_i^{\mathrm{dep}} - E_i^{\mathrm{arr}}}{\eta_c}
    \label{eq:energy_throughput}
\end{equation}

where the inner sum $\sum_{j \in \mathcal{E}_{n,t}} P^{\mathrm{c,max}}_{j}$ is the power already committed at node $n$ in period $t$ by the EVs assigned so far. This ensures that sufficient capacity remains at the node to fully charge EV $i$ within its active window, given the power already committed.

The committed load at each node is updated immediately after each successful assignment. Because~\eqref{eq:energy_throughput} ensures every assigned EV can be fully charged within its window, no SOC or node-capacity slack variables are needed in the LP.

\subsection{Battery Degradation Formulation}

V2G discharge accelerates battery wear. The degradation coefficient $\delta$ is approximated as:

\begin{equation}
\delta
=
\frac{\text{Battery Replacement Cost}}{\text{Lifetime Energy Throughput}}
\label{eq:deg_coeff}
\end{equation}

In this study a pack replacement cost of \$6{,}000 and a lifetime throughput of 300{,}000\,kWh (approximately 7{,}500 full equivalent cycles on a 40\,kWh pack) give $\delta = 0.02$\,\$/kWh.

The cost of degradation is a simplification of the underlying electrochemical processes. Whereas real degradation depends on depth-of-discharge, temperature, and charge rate, this study adopts a practical linear approximation for scheduling purposes, as it keeps the problem linear while still penalizing excessive discharge.

\subsection{Carbon Cap Formulation}

The uncoordinated emissions baseline assumes all EVs charge at maximum power on arrival, stopping once $E_{i,t}$ reaches the required departure energy:

\begin{equation}
\begin{aligned}
C^{\mathrm{base}}
&=
\sum_i \sum_{t \in \mathcal{T}^{\mathrm{base}}_i}
c_{n_i,t}\, P_i^{\mathrm{c,max}}
\end{aligned}
\label{eq:baseline}
\end{equation}

where $\mathcal{T}^{\mathrm{base}}_i \subseteq \mathcal{T}_i$ is the set of periods during which EV $i$ is still charging under the uncoordinated policy, that is, the periods before its required departure energy is reached. Restricting the sum to these periods is necessary; summing over the whole of $\mathcal{T}_i$ would overstate $C^{\mathrm{base}}$ by charging vehicles that are already full.

A soft carbon cap is imposed globally across the entire fleet, relative to this baseline:

\begin{equation}
\sum_i \sum_t
\left( P^{\mathrm{c}}_{i,t} - P^{\mathrm{d}}_{i,t} \right)
c_{n_i,t}
\;\le\;
\phi(\beta)\, C^{\mathrm{base}} + s^{\mathrm{carbon}}
\label{eq:carbon_cap}
\end{equation}

The slack variable $s^{\mathrm{carbon}} \ge 0$ intentionally softens the cap. Since carbon caps represent policy targets rather than hard physical limits, a nonnegative slack variable is included to ensure feasibility. A large penalty cost ($M = 10^6$) is assigned to the slack term in the objective, making cap violations highly expensive and therefore used only when necessary.

The carbon cap fraction ($\phi_0 = 0.85$) is a weight applied to the baseline emissions to determine the cap level. The effective cap fraction varies with the carbon preference weight $\beta$ as follows:

\begin{equation}
\phi(\beta) =
\begin{cases}
\text{disabled}           & \beta = 0 \\[4pt]
1 - \beta\,(1 - \phi_0)   & 0 < \beta < 1 \\[4pt]
\phi_0                    & \beta \ge 1
\end{cases}
\label{eq:phi_beta}
\end{equation}

where $\phi_0$ is the tightest allowable cap fraction. As $\beta$ increases, the carbon cap progressively tightens until it reaches $\phi_0$. When $\beta = 0$, the scheduler assigns no weight to emissions in the objective, so imposing a carbon cap would be contradictory, and the constraint is dropped entirely.

Coupling the cap to $\beta$ has a consequence that should be stated plainly. The cap is loosest where the objective produces the highest emissions and tightest where the objective has already driven emissions down, so it does not bind at any operating point reported in Section~\ref{sec:results}. It functions as a safeguard against schedules the objective would not select in any case, rather than as an active constraint. Tightening $\phi_0$ does not change this: the cap can only bind at high $\beta$, where emissions are already near the minimum attainable, so any $\phi_0$ low enough to bind there is infeasible and merely activates the slack term. Decoupling $\phi$ from $\beta$ and treating it as an exogenous policy limit would make the constraint active at the cost-dominated end of the frontier; this is left to future work.

\subsection{Optimization Problem}

\begin{subequations}
\label{eq:full_problem}
\begin{align}
\underset{
\substack{
P^{\mathrm{c}}_{i,t},
P^{\mathrm{d}}_{i,t},
E_{i,t},
s^{\mathrm{carbon}}
}
}{\text{minimize}}
\quad
&
\sum_{i \in \mathcal{E}}\; \sum_{t \in \mathcal{T}}
\Bigl[
  S_{n_i,t}
  \bigl( P^{\mathrm{c}}_{i,t} - P^{\mathrm{d}}_{i,t} \bigr)
  +
  \delta\, P^{\mathrm{d}}_{i,t}
\Bigr]
\;+\;
M\, s^{\mathrm{carbon}}
\label{eq:objective}
\\[6pt]
\text{s.t.}
\quad
& E_{i,\, t_i^{\mathrm{arr}}} = E_i^{\mathrm{arr}},
  \quad \forall\, i
\label{eq:init_soc}
\\[4pt]
& E_{i,t}
=
E_{i,t-1}
+
\left(
\eta_c P^{\mathrm{c}}_{i,t}
-
\frac{P^{\mathrm{d}}_{i,t}}{\eta_d}
\right),
\label{eq:battery_dynamics}
\\
& \qquad
\forall\, i,\;
t=t_i^{\mathrm{arr}}+1,\ldots,t_i^{\mathrm{dep}}
\nonumber
\\[4pt]
& E_{i,\, t_i^{\mathrm{dep}}} \ge E_i^{\mathrm{dep}},
  \quad \forall\, i
\label{eq:dep_soc}
\\[4pt]
& 0 \le E_{i,t} \le E_i^{\mathrm{cap}},
  \quad \forall\, i, t
\label{eq:battery_limit}
\\[4pt]
& 0 \le P^{\mathrm{c}}_{i,t} \le P_i^{\mathrm{c,max}},
  \quad \forall\, i, t
\label{eq:charge_limit}
\\[4pt]
& 0 \le P^{\mathrm{d}}_{i,t} \le P_i^{\mathrm{d,max}},
  \quad \forall\, i, t
\label{eq:discharge_limit}
\\[4pt]
& \sum_{i \in \mathcal{E}_n}
  \bigl( P^{\mathrm{c}}_{i,t} - P^{\mathrm{d}}_{i,t} \bigr)
  \le P_n^{\max},
  \quad \forall\, n, t
\label{eq:node_limit}
\\[4pt]
& \text{\eqref{eq:carbon_cap} (soft carbon cap)}
\nonumber
\\[4pt]
& s^{\mathrm{carbon}} \ge 0
\label{eq:slack_bound}
\end{align}
\end{subequations}

The objective in~\eqref{eq:objective} minimizes the ICS-weighted charging cost while penalizing battery degradation and any violation of the carbon cap through the slack variable $s^{\mathrm{carbon}}$. Since the charging allocation procedure guarantees feasible node assignments, all energy and power constraints remain hard. Because both the objective and constraints are linear, the resulting optimization problem is a convex LP, and the CBC solver returns a globally optimal solution.

All power quantities are expressed over a one-hour scheduling interval, so energy in kWh is numerically equal to power in kW and the interval length is omitted throughout for readability.

Equation~\eqref{eq:objective} is a scalarization rather than a monetary cost function, and its terms are not in a common unit: $S_{n_i,t}$ is dimensionless by construction, $\delta P^{\mathrm{d}}_{i,t}$ is in dollars, and $M s^{\mathrm{carbon}}$ is in grams. Combining them yields a dimensionless preference score whose minimizer is the schedule sought; the cost and emissions figures reported in Table~\ref{tab:baseline_comparison} are recomputed after the solve from the raw $p_{n,t}$ and $c_{n,t}$, not read off the objective value.

% =========================================================

\section{Results and Discussions} \label{sec:results}

\subsection{Simulation Setup}

\subsubsection{Data Collection}

LMP signals are taken from CAISO day-ahead hourly data for February 18, 2026, retrieved via the Open Access Same-Time Information System (OASIS) \cite{caiso_oasis}. Four nodes were selected for locational price diversity---CLAP\_BUNDLD-APND, POD\_DUTCH1\_7\_UNIT~1-APND, POD\_SLST13\_2\_SOLAR1-APND, and ALAMIT\_2\_PL1X3-APND---with profiles shown in Figure~\ref{fig:price}. They are placed at symmetrically spread positions on a synthetic grid, chosen for spatial diversity rather than to reflect true geography.

\begin{figure}[tbp]
\centering
\includegraphics[width=\linewidth]{fig_nodal_prices.pdf}
\caption{Nodal LMP across the four selected CAISO nodes on February 18, 2026, from CAISO OASIS.}
\label{fig:price}
\end{figure}

Carbon intensity is drawn from the CAISO average emission rate report for February 2026 \cite{caiso_carbon_intensity}, which reports a single California-wide average. Because no public source of locationally varying carbon intensity exists, each node is assigned a distinct daily profile: the 24-hour series for February 17 through 20 are mapped to the four nodes in the order listed above (Figure~\ref{fig:carbon}).

\begin{figure}[tbp]
\centering
\includegraphics[width=\linewidth]{fig_nodal_carbon.pdf}
\caption{Carbon intensity across February 17--20, 2026: day-specific CAISO average emission rates mapped to nodes as described in Section~\ref{sec:results}.}
\label{fig:carbon}
\end{figure}

A synthetic fleet of 30 EVs is used \cite{ev_fleet_data}. All vehicles share a 40\,kWh battery, 12\,kW maximum charging power, 4\,kW maximum discharging power, and 0.95 charge and discharge efficiency, and all must reach full charge before departure; they differ in initial SOC, arrival and departure time, and grid coordinates. Each node has an 85\,kW capacity, and the maximum assignment distance is 4.0 grid units. The fleet and node layout is shown in Figure~\ref{fig:grid_positions}.

\begin{figure}[tbp]
\centering
\includegraphics[width=0.7\linewidth]{grid_positions.pdf}
\caption{Starting locations of the 30-EV fleet and the four charging nodes.}
\label{fig:grid_positions}
\end{figure}

\subsubsection{Evaluation Schemes}

To evaluate the effectiveness of the scheduler, the following three schemes are tested alongside the uncoordinated baseline:

\begin{enumerate}

    \item \textbf{Cost-only optimization} ($\alpha = 1,\; \beta = 0$): optimized based solely on the LMP signal.

    \item \textbf{Carbon-only optimization} ($\alpha = 0,\; \beta = 1$): optimized based solely on the carbon intensity signal.

    \item \textbf{Balanced optimization} ($\alpha = 0.5,\; \beta = 0.5$): applies equal weight to LMP and carbon intensity.

\end{enumerate}

The uncoordinated baseline assigns each EV to its closest node and charges at maximum power on arrival until the required departure energy is reached. It provides the reference for all reported improvements. Schemes 1--3 all use the charging allocation procedure, V2G discharging, the carbon cap, and battery degradation, and correspond to the endpoints and midpoint of the trade-off frontier.

\subsubsection{Trade-off Sweep Setup}
\label{sec:sweep_setup}

To characterize the full cost--emissions trade-off, $\alpha$ is swept from 0 to 1 in uniform steps of 0.05 with $\beta = 1 - \alpha$, giving 21 operating points.

At every point the complete two-stage pipeline is re-run: the ICS in~\eqref{eq:composite_signal} is re-evaluated, the allocation of~\eqref{eq:assignment_score}--\eqref{eq:energy_throughput} is recomputed from scratch, and~\eqref{eq:full_problem} is re-solved against the resulting assignment. Node populations therefore differ between points, since a price-weighted run pulls vehicles toward the cheapest nodes and a carbon-weighted run toward the cleanest. Each point is the joint outcome of allocation and dispatch, not dispatch over a fixed allocation.

Total operating cost and emissions are recorded at each point and plotted as the frontier.

\subsubsection{Implementation}

The model is implemented in Python with PuLP and solved using CBC. Each instance covers 30 vehicles over 24 hourly periods: 1{,}440 charge and discharge variables, up to 720 state-of-charge variables, one slack variable, and roughly 2{,}400 constraints. A single instance solves to optimality in under a second, and the full study---baseline, three schemes, and the 21-point sweep---completes in under a minute on a standard workstation.

% ---------------------------------------------------------

\subsection{Overall Performance}
\label{sec:overall}

Table~\ref{tab:baseline_comparison} summarizes the total operating cost and carbon emissions produced by each scheduling scheme relative to the uncoordinated baseline.

\begin{table}[tbp]
\centering
\caption{Baseline vs.\ Optimized Charging: Cost and Emissions}
\label{tab:baseline_comparison}
\small
\setlength{\tabcolsep}{5pt}
\begin{tabular}{lcccccc}
\toprule
\textbf{Strategy}
  & \makecell{Total\\(\$)}
  & \makecell{Energy\\(\$)}
  & \makecell{Deg.\\(\$)}
  & \makecell{Savings\\(\%)}
  & \makecell{Emissions\\(g CO$_2$)}
  & \makecell{Reduction\\(\%)} \\
\midrule
Uncoordinated & 20.57 & 20.57 & 0.00 & ---  & 130,698.68 & ---  \\
Cost-only     & 13.67 & 11.83 & 1.84 & 33.6 & 115,465.27 & 11.7 \\
Carbon-only   & 19.26 & 16.86 & 2.40 &  6.4 &  98,670.32 & 24.5 \\
Balanced      & 14.74 & 12.42 & 2.32 & 28.4 & 102,494.53 & 21.6 \\
\bottomrule
\end{tabular}

\vspace{4pt}

{\footnotesize Total cost equals the sum of energy cost and battery degradation cost, as defined in~\eqref{eq:objective}. Savings and emission reductions are reported relative to the uncoordinated baseline.}
\end{table}

All three optimized schemes improve on both objectives. The carbon cap slack variable was never activated, and no fallback node assignments were required.

The cost-only scheme achieves the largest cost saving at 33.6\%, reducing total operating cost from \$20.57 to \$13.67, for a modest 11.7\% emissions reduction. The carbon-only scheme cuts emissions by 24.5\%, from 130{,}699\,g to 98{,}670\,g, but returns only 6.4\% on cost: it discharges most heavily of the three and so carries the highest degradation charge, \$2.40.

The balanced scheme performs well on both. It reduces cost by 28.4\% and emissions by 21.6\% at once, giving up 5.2 percentage points of cost saving against the cost-only scheme and 2.9 points of emissions reduction against the carbon-only scheme.

Cost is reported on the total-cost basis because energy cost alone is not what~\eqref{eq:objective} minimizes: the scheduler trades energy cost against battery throughput, so the energy-only figure is not guaranteed to fall monotonically in $\alpha$, while the total-cost figure is. The degradation penalty $\delta = 0.02$\,\$/kWh suppresses excessive cycling while still permitting profitable discharge when price or carbon signals are strong.

No binary variable prevents $P^{\mathrm{c}}_{i,t}$ and $P^{\mathrm{d}}_{i,t}$ from being simultaneously positive, and none is needed: with $\delta > 0$ and $\eta_c, \eta_d < 1$, doing both at once incurs a round-trip loss and a degradation charge for no benefit. Across all 720 EV-periods of the balanced schedule, no simultaneous charge--discharge occurs.

\subsection{Response to LMP Signals}

Figure~\ref{fig:cost_overlay} shows the per-node aggregate load overlaid with the LMP signal under the cost-only scheme. The scheduler charges when LMP is low and discharges via V2G when it is high, and this inverse relationship holds across all nodes. The few intervals that depart from it are driven by binding energy and node-capacity constraints, which push individual vehicles into less favorable periods while the fleet-wide schedule remains optimal.

\begin{figure}[tbp]
\centering
\includegraphics[width=\linewidth]{fig_cost_price_overlay.pdf}
\caption{Per-node aggregate fleet load compared with LMP signal under the cost-only scheme. Negative net power
indicates V2G discharge.}
\label{fig:cost_overlay}
\end{figure}

CLAP\_BUNDLD-APND serves nine EVs under this scheme and has the widest LMP spread, roughly \$0.01 to \$0.08/kWh. It shows the clearest load shifting, with charging concentrated in the low-price morning and a V2G discharge spike at the afternoon peak. At POD\_DUTCH1\_7\_UNIT~1-APND, where LMP turns negative at hours 7 and 9 and reaches \$\textminus0.025/kWh at hour 7, the largest charging power is observed.

\subsection{Response to Carbon Intensity}

Figure~\ref{fig:carbon_overlay} shows net load overlaid with carbon intensity under the carbon-only scheme. Charging concentrates in the midday low-carbon window (hours 9--16), when solar output drives intensity to its daily minimum.

Because each node carries its own emission profile, the depth and timing of that trough differ between nodes, and the scheduler shifts load accordingly rather than applying one fleet-wide pattern. Nodes with the deepest trough absorb the largest share of charging. This is the mechanism by which nodal, rather than system-wide, carbon intensity changes the schedule.

The evening ramp is equally informative. As solar falls and intensity climbs, several vehicles switch to V2G discharge, visible as negative net power. Discharging into a high-carbon period displaces generation that would otherwise come from fossil units, so the emissions credit exceeds the round-trip loss. This scheme discharges most heavily of the three and therefore carries the highest degradation charge in Table~\ref{tab:baseline_comparison}, \$2.40, which is the direct cost of the additional reduction it achieves.

\begin{figure}[tbp]
\centering
\includegraphics[width=\linewidth]{fig_carbon_overlay.pdf}
\caption{Aggregate fleet charging load compared with grid carbon intensity under the carbon-only scheme.  Negative net
power indicates V2G discharge.}
\label{fig:carbon_overlay}
\end{figure}

\subsection{Charging Allocation}
\label{sec:nodal}

Figure~\ref{fig:gridmap} shows the 30-EV fleet across the $10 \times 10$ grid and the resulting assignments under the balanced scheme. Because the allocation is ICS-driven, node populations depend on $(\alpha, \beta)$: the assignments shown are a single point of the sweep in Section~\ref{sec:pareto}, and every other value of $\alpha$ produces its own. The cost-only scheme, for instance, places nine EVs at CLAP\_BUNDLD-APND rather than the ten it receives under balanced weighting.

\begin{figure}[tbp]
\centering
\includegraphics[width=0.7\linewidth]{fig_grid_map.pdf}
\caption{Spatial visualization of EV-to-node assignments under the balanced scheme.  Dashed lines connect each EV to its assigned charging node; labels identify each EV by index.  The four CAISO nodes are shown as named anchors.}
\label{fig:gridmap}
\end{figure}

The balanced allocation is uneven: CLAP\_BUNDLD-APND serves 10 EVs, ALAMIT\_2\_PL1X3-APND 9, POD\_DUTCH1\_7\_UNIT~1-APND 6, and POD\_SLST13\_2\_SOLAR1-APND 5. Vehicles generally go to their nearest node, but not always: EV~10 is assigned to CLAP\_BUNDLD-APND rather than the closer POD\_DUTCH1\_7\_UNIT~1-APND. Both lie within $d_{\max}$ of that vehicle, so distance is not what separates them; ICS scoring and node capacity are. The cost of this reallocation is small, an 11\% increase in average nodal distance from 2.14 to 2.37 grid units.

\subsection{Cost--Emissions Trade-off Analysis}
\label{sec:pareto}

Figure~\ref{fig:pareto} shows the cost--emissions frontier obtained by sweeping $\alpha$ from 0 to 1 in steps of 0.05. Both allocation and dispatch are recomputed at each of the 21 points (Section~\ref{sec:sweep_setup}), so consecutive points may differ in where a vehicle charges as well as when.

One consequence is worth stating. Because Stage 1 is a greedy heuristic re-run at every point, the sweep is a parametric trace of the framework rather than a set that is non-dominated by construction, so non-dominance must be verified rather than assumed. Across the 21 points, total operating cost falls monotonically and emissions rise monotonically, so every value of $\alpha$ does yield a non-dominated solution on this instance. The frontier is mostly convex, steepening once the balanced point is passed.

Two regions stand out. Between $\alpha = 0$ and $\alpha = 0.05$, cost falls \$1.43 for only 177\,g more emissions, a large saving from even a small weight on price. Between $\alpha = 0.95$ and $\alpha = 1.00$, emissions rise 1{,}618\,g for a \$0.01 saving, the point at which further weight on price buys almost nothing. Both shifts come from the objective weighting; the carbon cap binds at neither point (Section~\ref{sec:overall}).

\begin{figure}[tbp]
\centering
\includegraphics[width=0.8\linewidth]{fig_pareto.pdf}
\caption{Trade-off frontier between total operating cost (energy plus battery
degradation) and total carbon emissions as $\alpha$ varies from 0 (carbon-only,
$\beta=1$) to 1 (cost-only, $\alpha=1$). The color gradient encodes $\alpha$.
The balanced operating point ($\alpha = \beta = 0.5$) is highlighted as the red
diamond.}
\label{fig:pareto}
\end{figure}

The balanced point ($\alpha = \beta = 0.5$, \$14.74, 102{,}495\,g) sits near the knee. Moving to the cost-only endpoint recovers \$1.07 (7.3\%) but adds 12{,}971\,g (12.7\%); moving to the carbon-only endpoint saves 3{,}824\,g (3.7\%) but costs \$4.52 (30.7\%). Emissions reduction slows sharply past the balanced point, which is what makes equal weighting an attractive operating choice.

\section{Conclusion}
\label{sec:conclusion}

This paper proposes a multi-objective EV charging and V2G discharging scheduling framework that minimizes both electricity cost and carbon emissions through an Integrated Charging Signal (ICS), which combines the LMP signal and grid carbon intensity. The scheduling framework consists of two stages: charging allocation and linear programming.

Three scheduling schemes were evaluated against the uncoordinated baseline, and all improved on both metrics. The cost-only scheme reduces total operating cost by 33.6\%, and the carbon-only scheme reduces emissions by 24.5\%, both relative to the baseline. The balanced scheme, which assigns equal weight to both signals, achieves a 28.4\% cost reduction alongside a 21.6\% emissions reduction. Each scheme thus achieves a large reduction in its target objective, and the balanced scheme improves substantially on both.

The trade-off analysis over $\alpha \in [0,1]$ (with $\beta = 1 - \alpha$) yields a mostly convex frontier of non-dominated solutions. The knee of the frontier lies close to the balanced point, making equal weighting an effective choice for both objectives.

% ---------------------------------------------------------

\subsection{Future Work}

Several simplifying assumptions were made in this study that may affect the results:

\begin{itemize}

\item Charging nodes were placed on a synthetic $10 \times 10$ grid rather than a real distribution feeder.

\item All charging nodes were assumed to have identical power capacities of $85\,\text{kW}$.

\item EV charging demand was assumed not to influence locational marginal prices.

\item Battery degradation was modeled as a linear function of discharged energy.

\item EV arrival times, departure times, and energy requirements were assumed to be known with certainty.

\item Travel distance was represented using Euclidean distance, neglecting road networks and traffic conditions that may influence charging-node selection.

\item Carbon intensity was represented using grid average values across several days rather than real nodal carbon intensity records.

\end{itemize}

Future work can address these limitations through several extensions:

\begin{itemize}

\item Integration with IEEE 33-bus or IEEE 69-bus distribution systems to incorporate realistic voltage and current constraints.

\item Development of stochastic or robust optimization frameworks that account for uncertainty in EV arrivals, departures, and charging demand.

\item Incorporation of electrochemical battery degradation models with depth-of-discharge-dependent aging behavior.

\item Use of transportation-network-based distance metrics and traffic-aware routing for more realistic charging allocation.

\item Integration of real-time or nodal carbon intensity signals.

\item Decoupling the carbon cap from the objective weighting so that it acts as an exogenous policy limit and binds at the cost-dominated end of the frontier.

\end{itemize}

% Flush any remaining floats before the bibliography so no figure lands
% after the references.

\FloatBarrier

\vspace{6pt}

\begin{adjustwidth}{-\extralength}{0cm}

\reftitle{References}

\begin{thebibliography}{999}

\bibitem{EV_adoption_growth}
International Energy Agency.
\textit{Global EV Outlook 2025}; IEA: Paris, France, 2025.
Available online: \url{https://www.iea.org/reports/global-ev-outlook-2025}
(accessed on 27 July 2026).

\bibitem{uncoordinated_charging_grid_impacts}
Nutkani, I.U.; Toole, H.; Fernando, N.; Loh, A.
Impact of EV charging on electrical distribution network and mitigating
solutions---A review.
\textit{IET Smart Grid} \textbf{2024}, \textit{7}, 485--502.

\bibitem{sortomme2011}
Sortomme, E.; Hindi, M.M.; MacPherson, S.D.J.; Venkata, S.S.
Coordinated charging of plug-in hybrid electric vehicles to minimize
distribution system losses.
\textit{IEEE Trans. Smart Grid} \textbf{2011}, \textit{2}, 198--205.

\bibitem{sundstrom2012}
Sundstr\"om, O.; Binding, C.
Flexible charging optimization for electric vehicles considering
distribution grid constraints.
\textit{IEEE Trans. Smart Grid} \textbf{2012}, \textit{3}, 26--37.

\bibitem{gan2013}
Gan, L.; Topcu, U.; Low, S.H.
Optimal decentralized protocol for electric vehicle charging.
\textit{IEEE Trans. Power Syst.} \textbf{2013}, \textit{28}, 940--951.

\bibitem{li2014dlmp}
Li, R.; Wu, Q.; Oren, S.S.
Distribution locational marginal pricing for optimal electric vehicle
charging management.
\textit{IEEE Trans. Power Syst.} \textbf{2014}, \textit{29}, 203--211.

\bibitem{huang2015dlmp}
Huang, S.; Wu, Q.; Oren, S.S.; Li, R.; Liu, Z.
Distribution locational marginal pricing through quadratic programming
for congestion management in distribution networks.
\textit{IEEE Trans. Power Syst.} \textbf{2015}, \textit{30}, 2170--2178.

\bibitem{sarker2016}
Sarker, M.R.; Dvorkin, Y.; Ortega-Vazquez, M.A.
Optimal participation of an electric vehicle aggregator in day-ahead
energy and reserve markets.
\textit{IEEE Trans. Power Syst.} \textbf{2016}, \textit{31}, 3506--3515.

\bibitem{hoehne2016}
Hoehne, C.G.; Chester, M.V.
Optimizing plug-in electric vehicle and vehicle-to-grid charge
scheduling to minimize carbon emissions.
\textit{Energy} \textbf{2016}, \textit{115}, 646--657.

\bibitem{xu2020ghg}
Xu, L.; Yilmaz, H.\"U.; Wang, Z.; Poganietz, W.-R.; Jochem, P.
Greenhouse gas emissions of electric vehicles in Europe considering
different charging strategies.
\textit{Transp. Res. D Transp. Environ.} \textbf{2020}, \textit{87}, 102534.

\bibitem{cheng2022carbon}
Cheng, K.-W.; Bian, Y.; Shi, Y.; Chen, Y.
Carbon-aware EV charging.
In Proceedings of the IEEE International Conference on Communications,
Control, and Computing Technologies for Smart Grids (SmartGridComm),
Singapore, 25--28 October 2022.

\bibitem{das2020}
Das, R.; Wang, Y.; Putrus, G.; Kotter, R.; Marzband, M.; Herteleer, B.;
Warmerdam, J.
Multi-objective techno-economic-environmental optimisation of electric
vehicle for energy services.
\textit{Appl. Energy} \textbf{2020}, \textit{257}, 113965.

\bibitem{aurangzeb2025}
Aurangzeb, M.; Wang, Y.; Iqbal, S.; Ullah, Z.
A novel multiobjective optimization approach for EV charging and
vehicle-to-grid scheduling strategy.
\textit{Int. Trans. Electr. Energy Syst.} \textbf{2025}, \textit{2025}, 1192925.

\bibitem{alikhani2025}
Alikhani, P.; Brinkel, N.; Schram, W.; Lampropoulos, I.; van Sark, W.
Multi-objective optimization of electric vehicle charging considering
market coupling in north-western Europe.
\textit{Transp. Res. D Transp. Environ.} \textbf{2025}, \textit{146}, 104829.

\bibitem{ambrosino2025}
Ambrosino, L.; Nguyen, K.M.; Vu, M.B.; Zorgati, R.; El Ghaoui, L.;
Calafiore, G.C.
A multi-criterion approach to smart EV charging with CO$_2$ emissions
and cost minimization.
\textit{arXiv} \textbf{2025}, arXiv:2511.06131.

\bibitem{moradipari2020}
Moradipari, A.; Alizadeh, M.
Pricing and routing mechanisms for differentiated services in an
electric vehicle public charging station network.
\textit{IEEE Trans. Smart Grid} \textbf{2020}, \textit{11}, 1489--1499.

\bibitem{arias2023lmp}
Arias, N.B.; Giraldo, J.S.; L\'{o}pez, J.C.; Hoogsteen, G.; Hurink, J.L.
EV allocation and charging within parking lots using a locational
marginal pricing mechanism.
In Proceedings of the IEEE PES Innovative Smart Grid Technologies Europe
(ISGT Europe), Grenoble, France, 10--13 October 2023; pp.\ 1--5.

\bibitem{caiso_oasis}
California Independent System Operator.
Open Access Same-Time Information System (OASIS): Day-Ahead Locational
Marginal Prices, trade date 18 February 2026.
Available online: \url{https://oasis.caiso.com}
(accessed on 27 July 2026).

\bibitem{caiso_carbon_intensity}
California Independent System Operator.
Average Emissions Rate Reports, CAISO Library.
Available online: \url{https://www.caiso.com/library/average-emissions-rate-reports}
(accessed on 27 July 2026).

\bibitem{ev_fleet_data}
Jung, J.
EV Fleet Simulation Data (ev\_infoV5.csv), GitHub repository, 2026.
Available online: \url{https://github.com/jiwanjung0101/EV-Charging-Project}
(accessed on 27 July 2026).

\end{thebibliography}

% \bibliography{references}

\end{adjustwidth}

\end{document}
